\FloatBarrier
\section{Implementation}
\label{sec:implementation}
\FloatBarrier
\subsection{High-Level Architecture}
Figure~\ref{fig:architecture} distinguishes TEPL's three inputs: a rule,
a dialect defining operations and their semantics, and a graph to rewrite.
Rules and dialects feed compilation, whereas the graph enters the runtime
engine directly. The implementation separates rule compilation from rewrite
execution. The additional-engine branch in the figure is an extension target,
not an implemented e-graph integration.


\begin{figuretext}{figures/TEPL architecture.pdf}{Rules and dialect semantics compile into reusable programs that a rewrite engine applies to input graphs.}{fig:architecture}{Rule and dialect cards feed the compiler and compiled JSON. A graph card
bypasses compilation and enters a greedy rewrite loop. A dashed adapter
connects compiled rules to a prospective additional engine.}
TEPL is implemented in C++20 and uses ANTLR to parse operation definitions
and rules. Compilation resolves names and checks types, constructs missing
attributes, checks semantic rules, and lowers completed rules into an
executable representation that can be exported as JSON:
\begin{codeblock}[H]
\begin{lstlisting}
Rule compilation
TEPL source
  -> ANTLR parsing
  -> Name resolution and type checking
  -> Attribute synthesis and semantic verification
  -> Executable rewrite program
  -> JSON export
\end{lstlisting}
\end{codeblock}
Type-checked rules may still contain holes; executable rules have completed
attribute programs, guards, and replacement patterns. JSON serializes this
executable representation for loading without repeating compilation or
semantic verification.

The rewrite engine loads and validates the program, searches the input graph
for an LHS match, evaluates its guard, computes replacement attributes, and
constructs and applies the RHS:
\end{figuretext}

\begin{codeblock}[H]
\begin{lstlisting}
Rewrite execution
Compiled JSON + input graph
  -> Load and validate
  -> Find an LHS pattern match
  -> Evaluate the guard
  -> Compute replacement attributes
  -> Construct and apply the RHS
\end{lstlisting}
\end{codeblock}
A failed guard or invalid attribute evaluation skips the match. The engine
rewrites greedily until no rule changes the graph or an execution limit is
reached. As described in Section~\ref{sec:auto}, generated Rust provides
another execution path with the shared runtime; C++ generation is unavailable.

\FloatBarrier
\subsection{Internal Representations}
Symbolic groups and position expressions use immutable typed nodes.
Types distinguish axes, positions, extents, integers, and Boolean conditions.
Gather, concatenation, and inverse permutation remain symbolic until their
inputs are known. Expression nodes can share subexpressions. Structural keys
identify expressions with the same structure, and each proof context interns
these keys to accelerate comparisons. Hashes support lookup, with structural
comparison resolving collisions.

Value nodes represent constants, operand reads, binary addition, binary
multiplication, and reduction. A separate index-expression representation
handles coordinate selection, concatenation, scaling, and offsets for
calculated reads. Normalization produces a vector of terms, each storing
an integer coefficient, a list of reads, and a symbolic reduction sequence,
as in Section~\ref{sec:normalization}.

\FloatBarrier
\subsection{Mapping Storage and Lookup}
Mapping entries store \code{(interface, group, positions)}. Separate indexes
support interface/group lookup and reverse lookup by interface and position
sequence. Partitions refer to entries through stable identifiers and retain
their own position sequences when these differ from the table's preferred
sequence. Overlap records retain the selected subgroup and positions in both
parent groups.

Lookup first checks existing mappings. If no direct mapping is available,
it follows established relationships and composes position sequences.
Concatenated groups combine component mappings. Traversal through registered
origins uses breadth-first search and avoids revisiting an origin. Results
are cached; registering mappings invalidates the mapping cache, and changes
to established facts invalidate the corresponding reasoning caches. Later
queries can therefore exploit newly discovered relationships.

\FloatBarrier
\subsection{Rule Execution}
The engine searches reachable graph nodes for LHS matches, checking operation
identity and operand structure and recording captured nodes, descriptors, and
shape bindings. Repeated tensor captures must refer to the same node.
The evaluator then processes the guard and replacement expressions using
those bindings. Boolean conditions preserve short-circuit evaluation; a
false guard or evaluation failure skips the match.

The builder reuses captured nodes and constructs new RHS operations with
the computed attributes. Construction occurs on a candidate graph, permitting
operation-contract and output-compatibility validation before committing
the replacement. After a successful replacement the engine restarts its
search. It stops when no candidate changes the graph, a configured limit
is reached, or execution is cancelled.

The mechanisms in this section describe the implementation status recorded
in the supplied system outline; they are separate from the intended universal
contract and from unperformed experiments. \todo{Pin the implementation and
runtime revisions and attach build, supported-operation, semantic-checker,
and generated-Rust validation evidence. The paper repository contains no
implementation artifact with which to independently audit these statements.}
